\documentclass[tikz,border=8pt]{standalone}
\usepackage{tikz}
\usepackage{amsmath}
\usepackage{amssymb}
\usepackage{ctex}
\usetikzlibrary{calc, positioning, arrows.meta, decorations.pathreplacing, fit, backgrounds}

\begin{document}
\begin{tikzpicture}[>=Stealth,
  box/.style={draw=gray!70, fill=white, rounded corners=5pt,
              minimum width=3.0cm, minimum height=1.4cm,
              inner sep=6pt, align=center, font=\small},
  subbox/.style={draw=#1!60, fill=#1!10, rounded corners=4pt,
                 minimum width=2.6cm, minimum height=1.1cm,
                 inner sep=5pt, align=center, font=\footnotesize},
  dombox/.style={draw=gray!50, fill=gray!5, rounded corners=6pt,
                 minimum width=3.8cm, minimum height=5.2cm},
]

%% ===== 左侧大框：R^n（定义域）=====
\node[dombox] (Rn) at (0,0) {};
\node[font=\small\bfseries, gray!70] at (0, 2.3) {$\mathbb{R}^n$};

% 行空间（Row Space）
\node[subbox=blue] (row) at (0, 1.0)
  {\textbf{行空间}\\$C(A^\top)$\\维数 $= r$};

% 零空间（Null Space）
\node[subbox=orange] (null) at (0, -1.0)
  {\textbf{零空间}\\$N(A)$\\维数 $= n-r$};

%% ===== 右侧大框：R^m（值域）=====
\node[dombox] (Rm) at (7,0) {};
\node[font=\small\bfseries, gray!70] at (7, 2.3) {$\mathbb{R}^m$};

% 列空间（Column Space）
\node[subbox=teal] (col) at (7, 1.0)
  {\textbf{列空间}\\$C(A)$\\维数 $= r$};

% 左零空间（Left Null Space）
\node[subbox=red] (lnull) at (7, -1.0)
  {\textbf{左零空间}\\$N(A^\top)$\\维数 $= m-r$};

%% ===== 矩阵 A 箭头 =====
\draw[->, very thick, gray!60] (1.95, 0.5) -- (5.05, 0.5)
  node[midway, above, font=\normalsize\bfseries] {$A$};
\draw[->, very thick, gray!60] (1.95, -0.5) -- (5.05, -0.5)
  node[midway, below, font=\small, gray] {$\mathbb{R}^n \to \mathbb{R}^m$};

%% ===== 正交补标注 =====
% 行空间 ⊥ 零空间（在 R^n 内）
\draw[<->, orange!70, thick, dashed] (0, 0.35) -- (0, -0.35)
  node[midway, right, font=\scriptsize, orange!80!black] {正交补};

% 列空间 ⊥ 左零空间（在 R^m 内）
\draw[<->, red!60, thick, dashed] (7, 0.35) -- (7, -0.35)
  node[midway, right, font=\scriptsize, red!80!black] {正交补};

%% ===== 同构箭头（行空间 ↔ 列空间）=====
\draw[->, blue!60, thick, bend left=20]
  (row.north east) to node[above, font=\scriptsize, blue!70!black]
  {$A\mathbf{x}$ 双射（维数同为 $r$）}
  (col.north west);

% 零空间映射到 0
\draw[->, orange!60, thick, bend right=20]
  (null.south east) to node[below, font=\scriptsize, orange!80!black]
  {$A\mathbf{x}=\mathbf{0}$（映到零向量）}
  (lnull.south west);

%% ===== 维数关系标注 =====
\node[font=\small, fill=white, draw=gray!50, thin, rounded corners=3pt, inner sep=6pt,
      align=center] at (3.5, -3.5) {
  \textbf{秩-零化度定理：} $\mathrm{rank}(A) + \mathrm{nullity}(A) = n$\\
  即：$r + (n-r) = n$，\quad $r + (m-r) = m$
};

%% ===== 顶部参数说明 =====
\node[font=\small, fill=white, draw=gray!50, thin, rounded corners=3pt, inner sep=5pt]
  at (3.5, 3.4) {矩阵 $A$ 为 $m\times n$ 阶，秩为 $r$};

%% 整体标题
\node[font=\large\bfseries] at (3.5, 4.4)
  {矩阵的四个基本子空间};

\end{tikzpicture}
\end{document}
